% Part: computability
% Chapter: recursive-functions
% Section: primitive-recursion

\documentclass[../../../include/open-logic-section]{subfiles}

\begin{document}

\olfileid{cmp}{rec}{pre}
\olsection{Rekursi Primitif}

Ciri wilangan asli yaiku saben wilangan asli bisa digayuh saka~$0$
kanthi ngetrapake operasi penerus~$+1$ kaping winates---saben
wilangan asli iku~$0$ utawa peneruse \dots{} peneruse~$0$.
Salah siji cara nemtokake fungsi~$h\colon\Nat \to \Nat$
kang migunakake kasunyatan iki kaya mangkene: (a)~tetepna
biji $h$ kanggo argumen~$0$, lan (b)~tetepna, yen biji
$h(x)$ wis dingerteni, carane ngitung biji $h(x+1)$.
Amarga (a)
langsung menehi biji $h(0)$, mula $h$~ditetepake kanggo~$0$.
Saiki, nganggo pituduh saka (b) kanggo $x=0$, kita bisa ngitung
$h(1) = h(0+1)$ saka~$h(0)$. Nganggo pituduh kang padha kanggo
$x=1$, kita ngitung $h(2) = h(1+1)$ saka~$h(1)$, lan sateruse.
% OLPL-121: Sumber mbalikke arah langkah pungkasan; aturan (b) lan
% tuladha ing ndhuwur nemtokake h(x+1) saka h(x), dudu kosokbaline.
Kanggo saben wilangan asli~$x$, pungkasane kita tekan langkah
kang netepake $h(x+1)$ saka $h(x)$, mula $h(x)$ ditetepake
kanggo kabeh $x \in \Nat$. Cathetan koreksi sumber OLPL-121: ukara sumber mbalikke arah langkah pungkasan; aturan rekursi kang dituduhake nemtokake nilai sabanjure saka nilai sadurunge.

Umpamane, upamana kita nemtokake $h\colon \Nat \to \Nat$
nganggo rong persamaan iki:
\begin{align*}
h(0) & =  1\\
h(x+1) & =  2 \cdot h(x)
\end{align*}
Yen kita wis ngerti carane ngalikake, persamaan iki menehi
katrangan kang dibutuhake kanggo (a) lan~(b) ing ndhuwur.
Kanthi bola-bali ngetrapake persamaan kapindho, kita entuk
\begin{align*}
  h(1) & = 2\cdot h(0) = 2,\\
  h(2) & = 2\cdot h(1) = 2\cdot 2,\\
  h(3) & = 2 \cdot h(2) = 2\cdot 2 \cdot 2,\\
  & \vdots
\end{align*}
Kita weruh yen fungsi~$h$ kang wis ditemtokake yaiku $h(x) = 2^x$.

Ciri khas wilangan asli njamin yen mung ana siji fungsi~$h$
kang nyukupi rong syarat iki. Sepasang persamaan kaya mangkene
diarani \emph{definisi kanthi rekursi primitif} kanggo
fungsi~$h$. Iki diarani “rekursif,” yaiku kita netepake~$h$
kanthi definisi kang, mligine ing persamaan kapindho,
ngemot $h$ dhewe ing sisih tengen. Iki diarani “primitif”
amarga kanggo netepake $h(x+1)$ kita mung nggunakake
biji~$h(x)$, yaiku biji sakdurunge kang langsung.
Iki cara paling prasaja kanggo netepake fungsi ing~$\Nat$
kanthi rekursif.

Kita bisa netepake fungsi kang luwih dhasar maneh, kayata
panjumlahan lan perkalian, kanthi rekursi primitif. Nanging
ing kene fungsi kang dirembug nduweni $2$ panggonan.
Kita netepake salah siji panggonan argumen, lan nggunakake
liyane kanggo rekursi. Umpamane, kanggo netepake
$\Add(x, y)$, kita bisa netepake~$x$ lan nemtokake biji
dhisik kanggo $y=0$, banjur kanggo $y+1$ adhedhasar~$y$.
Amarga $x$ tetep, argumen kasebut bakal katon ing sisih kiwa lan tengen
persamaan definisine.
\begin{align*}
\Add(x,0) & =  x\\
\Add(x,y+1) & =  \Add(x,y)+1
\end{align*}
Persamaan iki nemtokake biji $\Add$ kanggo kabeh $x$
\emph{lan}~$y$. Kanggo nemokake $\Add(2,3)$, umpamane,
kita ngetrapake persamaan definisi kanggo $x = 2$,
nganggo persamaan kapisan kanggo nemokake
$\Add(2,0) = 2$, banjur nganggo persamaan kapindho
kanthi mbaka siji nemokake $\Add(2,1) = 2 + 1 = 3$,
$\Add(2, 2) = 3 + 1 = 4$, $\Add(2,
3) = 4 + 1 = 5$.

Ing definisi $\Add$ kita nggunakake $+$ ing sisih tengen
persamaan kapindho, nanging mung kanggo nambah~$1$.
Tegese, kita nggunakake fungsi penerus $\Succ(z) = z+1$
lan ngetrapake marang biji sadurunge $\Add(x,y)$
kanggo netepake $\Add(x,y+1)$. Mula kita bisa nganggep
definisi rekursif iku diwenehi liwat siji fungsi kang
ditrapake marang biji sadurunge. Nanging ora apa-apa---lan
kadhangkala perlu---ngidini fungsi kasebut gumantung ora
mung marang biji sadurunge, nanging uga marang $x$ lan~$y$.
Coba delengen:
\begin{align*}
  \Mult(x,0) & =  0 \\
  \Mult(x,y+1) & =  \Add(\Mult(x,y),x)
\end{align*}
Iki definisi rekursif primitif fungsi $\Mult$ kanthi
ngetrapake fungsi $\Add$ marang biji sadurunge
$\Mult(x,y)$ lan argumen kapisan~$x$. Iki uga nemtokake
fungsi~$\Mult(x,y)$ kanggo saben argumen $x$ lan~$y$.
Umpamane, $\Mult(2,3)$ ditemtokake kanthi mbaka siji
ngitung $\Mult(2,0)$, $\Mult(2,1)$, $\Mult(2,2)$,
lan~$\Mult(2,3)$:
\begin{align*}
  \Mult(2,0) & = 0\\
  \Mult(2,1) & = \Mult(2,0+1) =
  \Add(\Mult(2,0), 2) = \Add(0, 2) = 2\\
  \Mult(2,2) & = \Mult(2,1+1) =
  \Add(\Mult(2,1), 2) = \Add(2, 2) = 4\\
  \Mult(2,3) & = \Mult(2,2+1) =
  \Add(\Mult(2,2), 2) = \Add(4, 2) = 6
\end{align*}

Mula pola umume kaya mangkene: kanggo menehi definisi rekursif
primitif fungsi~$h(x_0, \dots, x_{k-1}, y)$, kita nyedhiyakake
rong persamaan. Persamaan kapisan netepake biji
$h(x_0, \dots, x_{k-1}, 0)$ tanpa ngrujuk marang~$h$.
Persamaan kapindho netepake biji $h(x_0,
\dots, x_{k-1}, y+1)$ adhedhasar $h(x_0, \dots, x_{k-1}, y)$,
argumen liyane $x_0$, \dots,~$x_{k-1}$, lan~$y$.
Mung biji sadurunge saka~$h$ kang langsung bisa digunakake
ing persamaan kapindho. Yen operasi ing sisih tengen rong
persamaan kasebut dianggep minangka fungsi $f$ lan~$g$,
mula pola umum kanggo netepake fungsi anyar~$h$ kanthi
rekursi primitif kaya mangkene:
\begin{align*}
  h(x_0, \dots, x_{k-1}, 0) & = f(x_0, \dots, x_{k-1})\\
  h(x_0, \dots, x_{k-1}, y+1) & =
  g(x_0, \dots, x_{k-1}, y, h(x_0, \dots, x_{k-1}, y))
\end{align*}
Ing kasus $\Add$, kita nduweni $k=1$ lan $f(x_0) = x_0$
(fungsi identitas), lan $g(x_0, y, z) = z + 1$
(fungsi $3$ panggonan kang ngasilake penerus saka argumen
katelu):
\begin{align*}
  \Add(x_0, 0) & = f(x_0) = x_0\\
  \Add(x_0, y+1) & = g(x_0, y, \Add(x_0, y)) =
  \Succ(\Add(x_0, y))
\end{align*}
Ing kasus $\Mult$, kita nduweni $f(x_0) = 0$
(fungsi konstan kang tansah ngasilake~$0$) lan
$g(x_0, y, z) = \Add(z,x_0)$ (fungsi $3$ panggonan
kang ngasilake jumlah argumen pungkasan lan kapisan):
\begin{align*}
  \Mult(x_0,0) & =  f(x_0) = 0 \\
  \Mult(x_0,y+1) & =  g(x_0, y, \Mult(x_0,y)) =
  \Add(\Mult(x_0,y), x_0)
\end{align*}

\end{document}
